\section{Methodology}

The system contains four layers, summarized in Figure~\ref{fig:method-pipeline}. First, a legal synthetic generator creates fragments conditioned on pitch-class sets, optional interval vectors, twelve-tone rows, P/R/I/RI forms, rhythmic profiles, gesture labels, instrument, voice count, and measure count. Second, theory utilities compute pc-set coverage, interval vectors, row-order accuracy, aggregate completion, density curves, gesture features, and range violations. Third, an event tokenizer represents conditions and score events as Transformer-readable sequences. Fourth, generation uses a rule baseline or neural model with non-differentiable constraint penalties applied during decoding or reranking.

All outputs remain symbolic. MusicXML export is used for notation inspection, and analysis reports are saved alongside generated scores.

The proposed decoding procedure samples candidate symbolic continuations and reranks them with explicit post-tonal penalties. These penalties are not treated as differentiable training losses. They are used only where the metric is naturally symbolic, such as pc-set coverage, aggregate completion, row-order accuracy, rhythmic-profile distance, gesture consistency, and instrument range validity.
